\section{Arytmetyka końców Hilberta} % lang-pl
\label{sekcja_ciala_koncow}

Z samej geometrii płaszczyzny hiperbolicznej, bez żadnych liczb, Hilbert wydobędzie ciało uporządkowane -- i to jest jeden z najpiękniejszych pomysłów całej teorii. % lang-pl
W tej sekcji zobaczymy, jak z końców prostych zrobić ciało, jak odczytać w nim ruchy sztywne i kąty oraz jak dzięki niemu udowodnić dwa twierdzenia: o konstrukcji Bolyaia i o wysokościach trójkąta. % lang-pl

Hartshorne \cite[s. 388]{hartshorne_2000} opisuje tę konstrukcję jako pokaz siły abstrakcyjnej algebry -- tak jak w geometrii euklidesowej ciało arytmetyki odcinków pomagało zrozumieć płaszczyznę, tak tu ciało końców tłumaczy geometrię hiperboliczną; po przejściu do niego geometria „przeszła przemianę w coś bogatego i dziwnego''. % lang-pl

Konstrukcja zaczyna się od trzech prostych w dowolnym trójkącie: symetralne boków trójkąta w płaszczyźnie hiperbolicznej albo przecinają się w punkcie, albo mają wspólną prostopadłą, albo są równoległe graniczne o wspólnym końcu \cite[s. 388--389]{hartshorne_2000} (prop. 41.1; zobacz też sekcję \ref{sekcja_punkty_idealne}). % lang-pl
Z tego wyprowadza się twierdzenie o trzech odbiciach: jeśli proste $a, b, c$ mają wspólny koniec $\omega$, to istnieje czwarta prosta $d$ o tym samym końcu, że $\sigma_c\sigma_b\sigma_a = \sigma_d$ \cite[s. 389]{hartshorne_2000} (prop. 41.2). % lang-pl

Ustalmy prostą i oznaczmy jej końce przez $0$ i $\infty$; niech $F$ będzie zbiorem wszystkich końców różnych od $\infty$. % lang-pl
Dla dwóch końców $\alpha, \beta$ ich \emph{sumą} $\alpha + \beta$ jest koniec symetralnej odcinka $AB$ (różny od $\infty$), gdzie $A$ i $B$ powstają z punktu $C$ prostej $(0,\infty)$ przez odbicia w prostych $(\alpha,\infty)$ i $(\beta,\infty)$ \cite[s. 390]{hartshorne_2000}. % lang-pl
\emph{Iloczyn} definiuje się podobnie, po ustaleniu prostopadłej do $(0,\infty)$ i jej końców $1$ i $-1$, przez dodawanie odcinków ze znakiem \cite[s. 391]{hartshorne_2000}. % lang-pl

\begin{theorem}
	Zbiór $F$ końców z tak określonymi działaniami $+, \cdot$ i z pojęciem elementów dodatnich (końców leżących po tej samej stronie prostej $(0,\infty)$ co $1$) jest ciałem uporządkowanym euklidesowym. % lang-pl
\end{theorem}

Dowód (rozłożony na kolejne własności; rozdzielność korzysta z przesunięcia wzdłuż prostej): Hartshorne \cite[s. 390--392]{hartshorne_2000} (prop. 41.3, tw. 41.4). % lang-pl
\index{ciało!końców} % lang-pl

W tym ciele ruchy sztywne mają prostą postać: odbicie w $(0,\infty)$ to $x' = -x$, odbicie w $(1,-1)$ to $x' = 1/x$, przesunięcie wzdłuż $(0,\infty)$ to $x' = ax$, obrót wokół $\infty$ to $x' = x + a$, a obrót wokół początku układu to $x' = (x+a)/(-ax+1)$ \cite[s. 392--393]{hartshorne_2000} (prop. 41.5). % lang-pl
Prosta to nieuporządkowana para różnych końców $(u_1, u_2)$, a punkt -- zbiór wszystkich prostych przez niego przechodzących, opisany równaniem $u_1u_2 - b(u_1+u_2) + a^2 = 0$ z $a > 0$, $|b| < a$ \cite[s. 394]{hartshorne_2000} (prop. 41.6): jest to odwrócenie zwykłej geometrii analitycznej, w której punkt ma współrzędne, a prosta równanie. % lang-pl

Odcinek $AB$ ma \emph{multiplikatywną długość} $\mu(AB) = a$ (kładziemy go na promień $0\infty$, a prosta prostopadła w jego końcu to $(a, -a)$), która ma cztery naturalne własności: jest $> 1$, przystające odcinki mają równe $\mu$, $\mu$ rośnie z długością, a przy dodawaniu odcinków się mnoży \cite[s. 394--395]{hartshorne_2000} (prop. 41.7). % lang-pl
Podobnie dla kąta $\theta$ definiuje się $\tan\tfrac{\theta}{2} = a$ i wzory na sumę \cite[s. 395--396]{hartshorne_2000} (prop. 41.8); pozwala to zapisać wzór Bolyaia $\tan\tfrac{\alpha}{2} = \mu(AB)^{-1}$ \cite[s. 396]{hartshorne_2000} (prop. 41.9) i kąt między prostymi $(u_1,u_2)$, $(v_1,v_2)$ przez dwustosunek końców: $\tan\tfrac{\theta}{2} = \sqrt{-(v_1,v_2;u_1,u_2)}$ \cite[s. 399]{hartshorne_2000} (prop. 41.12). % lang-pl
Dwustosunek pojawił się już u Cayleya i Kleina (zobacz sekcję \ref{sekcja_beltrami_cayley_klein}); tu wynika z geometrii, a nie ją definiuje. % lang-pl

\begin{theorem}
[konstrukcja równoległej granicznej Bolyaia] % lang-pl
\index{konstrukcja!Bolyaia (równoległej granicznej)} % lang-pl
	Niech $PQ$ będzie prostopadłą z $P$ na prostą $l$, $m$ prostą przez $P$ prostopadłą do $PQ$, $R$ dowolnym punktem na $l$, a $RS$ prostopadłą z $R$ na $m$. % lang-pl
	Wtedy okrąg o środku $P$ i promieniu $QR$ przecina odcinek $RS$ w punkcie $T$, a półprosta $PT$ jest półprostą równoległą graniczną do $l$ przez $P$. % lang-pl
\end{theorem}

Dowód (rachunek w ciele końców): Hartshorne \cite[s. 396--398]{hartshorne_2000} (prop. 41.10, lemat 41.11). % lang-pl
Hartshorne podkreśla, że konstrukcja wymaga cyrkla i linijki, a jej działanie dowodzi się, zakładając wcześniej (z aksjomatu (L)), że równoległa graniczna istnieje; bez tego założenia wynik może zawieść -- na przykład w płaszczyznach nieArchimedesowych \cite[s. 398]{hartshorne_2000}. % lang-pl
Cytuje też Greenberga, według którego znalezienie bezpośredniego geometrycznego dowodu tej konstrukcji byłoby warte natychmiastowego doktoratu \cite[s. 398]{hartshorne_2000}. % lang-pl
Greenberg \cite[s. 210]{greenberg_2010} zauważa, że konstrukcja Bolyaia działa w każdej płaszczyźnie Hilberta spełniającej hipotezę kąta ostrego i aksjomat prosta--okrąg, ale półprosta nie musi być wtedy równoległa graniczna, jeśli nie ma dodatkowego założenia. % lang-pl

\begin{theorem}
[o wysokościach trójkąta] % lang-pl
	Jeśli dwie wysokości trójkąta w płaszczyźnie hiperbolicznej się przecinają, to trzecia też przechodzi przez ten punkt; jeśli dwie są równoległe graniczne, to wszystkie trzy mają wspólny koniec; jeśli dwie mają wspólną prostopadłą, to wszystkie trzy mają tę samą wspólną prostopadłą. % lang-pl
\end{theorem}

Dowód: Hartshorne \cite[s. 399--400]{hartshorne_2000} (prop. 41.13); zad. 40.13--40.14 podają dowód czysto geometryczny \cite[s. 387]{hartshorne_2000}. % lang-pl

Zadania (grupa ruchów sztywnych jako grupa przekształceń ułamkowo-liniowych, znajdowanie końców dwusiecznej itd.): Hartshorne \cite[s. 400--402]{hartshorne_2000}. % lang-pl
% Źródła: Hartshorne s. 403--414 (§42; zad. 42.1--42.25); Eves nie omawia trygonometrii.

